\section{Introduction}

\subsection{Motivation and Related Work}

\begin{itemize}
    \item Graph colouring is a fundamental problem in computer science, NP-hard for \(k\ge 3\) \cite{MR4679190-Karp72}.
    \item \whp random recovery is easy for \(k\)-colorings in random graphs \cite{MR1344544-Blum95, MR1484153-Alon97} though
    \item Key idea is to use the spectrum of the graph, and recover the colouring from the extreme eigenvectors. For this we need to guarantee they don't ``hide'' in the spectrum.
    \item Goal is pushing random as far as possible remove as many assumptions as possible.
    \item Semi-random setting: David and Feige~\cite{MR3536556-David16}, then derandomised by Kumar, Louis, and Tulsiani~\cite{kumar2018finding}.
    \item Related setting where one large independet set is to be found: Bafna, Hsieh, and Kothari~\cite{bafna-hsieh-kothari}. Mention the $k=3$ vs $k>3$ difference
    \item Explain Sum of Squares idea briefly and how it relates to recovery problems in general
    \item Another line of work is to use as few colours as possible (we don't consider this here).
\end{itemize}

\subsection{Results}


\begin{itemize}
    \item
\end{itemize}

\subsection{Techniques}

\begin{itemize}
    \item As demonstration give the simple example of linear degree semi-random
\end{itemize}

\subsection{EXTRA RESULTS TO ADD}

\begin{itemize}
    \item For $k=3$ we can use SoS to get low edge violation and show that it implies low vertex violation.
\end{itemize}

% [Introduction to coloring. Coloring random graphs: Blum and Spencer~\cite{MR1344544-Blum95}, Alon and Kahale~\cite{MR1484153-Alon97}.]

Finding proper \(k\)-colorings in graphs for \(k\ge 3\) is among the first computational problems shown to be NP-hard \cite{MR4679190-Karp72}.
% TODO; discuss approximation algorithms for coloring
As a consequence, much research has focused on approximation algorithms for this problem and identifying assumptions on the input that allow us to solve the problem efficiently \cite{MR819128-83, MR1369207-94, MR1344544-Blum95, MR1439867-97, MR1484153-Alon97, MR1623197-98, MR1774844-00, MR1948749-01, MR2277147-06, MR2538841-09, MR2863244-Arora11, MR3536556-David16, MR3634492-17, kumar2018finding, MR4764817-24}.

Most relevant to this paper, Blum and Spencer~\cite{MR1344544-Blum95} and Alon and Kahale~\cite{MR1484153-Alon97} studied an average-case setting in which a $k$-coloring is planted in a random graph,
and showed that polynomial-time algorithms can recover the planted $k$-coloring with high probability.
This result was subsequently derandomized by David and Feige~\cite{MR3536556-David16}:
They obtain the same guarantees if the $k$-coloring is planted \emph{at random} in a (two-sided) expander graph, where we require all but one eigenvalue of its normalized adjacency matrix to be close to \(0\).
Furthermore, they can recover a $k$-coloring for a $0.99$ fraction of the vertices even if the $k$-coloring is planted \emph{adversarially} in such a graph.

Kuman, Louis, and Tulsiani~\cite{kumar2018finding} provide a different derandomization result:
given a graph that admits a $3$-coloring satisfying a particular pseudo-randomness condition, their algorithms finds a $3$-coloring for a $0.99$ fraction of the vertices in time $n^{O(t)}$, where $t$ is the number of eigenvalues of the graph smaller than $-0.51$.

\paragraph{Two-sided vs one-sided expansion.}

The results of Alon-Kahale and David-Feige crucially exploit the fact that both the positive and negative part of the spectrum of the graph are severely restricted.

% [More recent, derandomization: Arora and Ge~\cite{MR2863244-Arora11}, David and Feige~\cite{MR3536556-David16}, Kumar, Louis, and Tulsiani~\cite{kumar2018finding}.]

% \rares{This paragraph repeats information from before}
% More recently, David and Feige~\cite{MR3536556-David16} derandomized various aspects of~\cite{MR1484153-Alon97}.
% Of most relevance to us, they obtain the same guarantees even if the balanced $k$-coloring is planted at random in a regular two-sided expander graph, instead of a random regular graph.
% Here, we call a graph is a two-sided expander if $\max\{\lambda_2, -\lambda_n\}$ is bounded away from $1$, where $1 = \lambda_1 \geq ... \geq \lambda_n$ are the eigenvalues of the normalized adjacency matrix of the graph.
% Furthermore, the algorithms of David and Feige also find a $k$-coloring on $0.99n$ of the vertices of the graph even if the $k$-coloring is planted adversarially instead of randomly.

Recent seminal work developed the first algorithms for one-sided expander graphs, where the negative part of the spectrum is completely unrestricted \cite{bafna-hsieh-kothari}.
The current work further develops our understanding of the computational landscape of coloring and independent set problems on one-sided expander graphs and provides answers to key questions left open by this prior work.  

% [Previous works use two-sided expanders, or more generally small bottom threshold rank. One-sided expanders: Bafna, Hsieh, and Kothari~\cite{bafna-hsieh-kothari}. They can find small independent sets.]

% \paragraph{Bottom threshold rank}

% [Explanation...]

% [Our contribution]

% [(1) One-sided expanders have constant threshold rank. Also generalization to bounded top threshold rank. State result. Say it is technically intricate.]

% [This does not trivialize problem. As~\cite{bafna-hsieh-kothari} showed, it is still hard to even find small independent sets in 4-colorable/6-colorable one-sided expanders.]

% [(2) Need pseudo-randomness in coloring. Indeed, we can recover a 0.99 k-coloring for any pseudo-random k-colorable one-sided expander. (To figure out if we want to discuss non-balancedness here.) If random planted, as in~\cite{MR3536556-David16}, can even get full coloring.]

% [(3) However, 3-colorable one-sided expanders are always pseudo-random. Therefore, we can recover a 0.99 3-coloring for any pseudo-random 3-colorable one-sided expander.]

% [(4) Finally,~\cite{bafna-hsieh-kothari} have a result for graphs with $(1/2-\e)$-independent set. We can get here a $(1/2-\e')$-independent set.]

% [Optional: to figure out whether/where to state hardness results.]




% \subsection{Notation}
% \label{sec:notation}

% \paragraph{General notations} For positive integers $l\leq l'$ , let $[l]$ denote $\Set{1\ldots, l}$ and let $\btw{l}{l'}$ denote $\Set{l\ldots, l'}$.

% \paragraph{Graph notations}
% Given a graph $G=(V, E)$, we denote by $A_G$ its unnormalized adjacency matrix, and by $\bar{A}_G = \frac{1}{d_{\avg}(G)} A_G$ its normalized adjacency matrix (we omit the subscript when it is clear from context).
% \david{This notion of normalized adjacency matrix is non standard and potentially problematic.}

% Given disjoint vertex sets $S, T \subseteq V$, we use $E(S, T)$ to denote the set of edges with one endpoint in $S$ and one endpoint in $T$. When $S = \set{s}$ is a single vertex, we simplify the notation to $E(s, T)$ (analogously for the case when $T$ is a single vertex).

% \paragraph{Matrix notations} Given matrix $M$, we denote by $\lambda_i(M)$ the $i$-th largest eigenvalue of $M$.

% \subsection{Organization}
% \label{sec:organization}



%%% Local Variables:
%%% mode: LaTeX
%%% TeX-master: "../main.tex"
%%% End:
